\setcounter{dang}{0}

\section{DÃY SỐ}
\subsection{LÝ THUYẾT CẦN NHỚ}
\subsubsection{ĐỊNH NGHĨA DÃY SỐ}
\iconMT \indam{ Định nghĩa dãy số:} Mỗi hàm số $u$ xác định trên tập các số nguyên dương $\mathbb{N}^*$ được gọi là một dãy số vô hạn (gọi tắt là dãy số). Kí hiệu $u=u(n )$.
	% \khung{
		\begin{listEX}[1]
			\item [\ding{172}] Ta thường viết $u_n$ thay cho $u(n)$ và kí hiệu dãy số $u=u(n)$ bởi $\left(u_n\right)$. Do đó dãy số $(u_n)$ được viết dưới dạng khai triển ${{u}_{1}},{{u}_{2}},{{u}_{3}},\ldots,{{u}_n},\ldots,$ trong đó
			\begin{itemize}
				\item [$\bullet$]  ${{u}_{1}}$ là số hạng đầu;
				\item [$\bullet$]  ${{u}_n}$ là số hạng thứ $n$ và là số hạng tổng quát của dãy số.
			\end{itemize}
			\item [\ding{173}] Nếu $\forall n \in \mathbb{N}^*, u_n=c$ thì $\left(u_n\right)$ được  gọi là dãy số không đổi.
		\end{listEX}
	% }
\iconMT \indam{Định nghĩa dãy số hữu hạn:}
	\begin{itemize}
		\item [$\bullet$] Mỗi hàm số $u$ xác định trên tập $M=\left\{ 1,2,3,\ldots,m \right\}$ với $m\in \mathbb{N}^*$ được gọi là một dãy số hữu hạn.
		\item [$\bullet$]  Dạng khai triển của nó là ${{u}_{1}},{{u}_{2}},{{u}_{3}},\ldots,{{u}_m},$ trong đó ${{u}_{1}}$ là số hạng đầu, ${{u}_m}$ là số hạng cuối.
	\end{itemize}

\subsubsection{CÁCH CHO MỘT DÃY SỐ }
Ta thường gặp một trong các cách sau đây:
\begin{listEX}[1]
	\item [\ding{172}] Liệt kê các số hạng (chỉ dùng cho các dãy hữu hạn và có ít số hạng);
	\item [\ding{173}] Dãy số cho bằng công thức của số hạng tổng quát;
	\item [\ding{174}] Dãy số cho bằng phương pháp mô tả;
	\item [\ding{175}] Dãy số cho bằng phương pháp truy hồi, nghĩa là
		\begin{itemize}
			\item [$\bullet$] Cho số hạng đầu (hay vài số hạng đầu).
			\item [$\bullet$] Cho hệ thức truy hồi, tức là hệ thức biểu thị số hạng thứ $n$ qua số hạng (hay vài số hạng) đứng trước nó.
		\end{itemize}
\end{listEX}

\subsubsection{DÃY SỐ TĂNG, DÃY SỐ GIẢM}
\begin{itemize}
\item [\ding{172}] Dãy số $({{u}_n})$ được gọi là dãy số tăng nếu ta có ${{u}_{n+1}}>{{u}_n}$ với mọi $n\in \mathbb{N}^*$.
\item  [\ding{173}] Dãy số $({{u}_n} )$ được gọi là dãy số giảm nếu ta có ${{u}_{n+1}}<{{u}_n}$ với mọi $n\in \mathbb{N}^*$.
\end{itemize}

\subsubsection{DÃY SỐ BỊ CHẶN}

\begin{itemize}
\item [\ding{172}] Dãy số $({{u}_n} )$ được gọi là bị chặn trên nếu tồn tại một số $M$ sao cho ${{u}_n}\le M,\forall n\in \mathbb{N}^*$.
\item [\ding{173}] Dãy số $({{u}_n} )$ được gọi là bị chặn dưới nếu tồn tại một số $m$ sao cho
${{u}_n}\ge m,\forall n\in \mathbb{N}^*$.
\item [\ding{174}] Dãy số $({{u}_n} )$ được gọi là bị chặn nếu nó vừa bị chặn trên vừa bị chặn dưới, tức là tồn tại các số $m$, $M$ sao cho
$m\le {{u}_n}\le M,\forall n\in \mathbb{N}^*$.
\end{itemize}
\subsection{PHÂN LOẠI VÀ PHƯƠNG PHÁP GIẢI TOÁN}
\begin{dang}{Tìm các số hạng của dãy số cho bởi công thức tổng quát}
\end{dang}
\begin{vd}
	Cho dãy số $\left(u_n\right)$, biết $u_n=\dfrac{n}{3^n-1}$. Tìm ba số hạng đầu tiên của dãy số.
	\dapso{$\dfrac{1}{2};\dfrac{1}{4};\dfrac{3}{26}$.}
	\loigiai{
		Ta có $u_1=\dfrac{1}{2};u_2=\dfrac{2}{3^2-1}=\dfrac{2}{8}=\dfrac{1}{4};u_3=\dfrac{3}{3^3-1}=\dfrac{3}{26}.$
	}
\end{vd}\dongcham{5}

\begin{vd}%[Đỗ Vũ Minh Thắng]%[1D3B2]
	Cho dãy số $(u_n)$ được xác định bởi $u_n=\dfrac{n^2+3n+7}{n+1}$.
	\begin{tasks}(1)
		\task Viết năm số hạng đầu của dãy. \dapso{$\dfrac{11}{2}$; $\dfrac{17}{3}$; $\dfrac{25}{4}$; $7$; $\dfrac{47}{6}$.}
		\task Dãy số có bao nhiêu số hạng nhận giá trị nguyên?
		\dapso{$ 1 $.}
	\end{tasks}
	
	\loigiai{
		
		\begin{enumerate}
			\item Ta có năm số hạng đầu của dãy
			$u_1=\dfrac{1^2+3.1+7}{1+1}=\dfrac{11}{2}$; $u_2=\dfrac{17}{3}$; $u_3=\dfrac{25}{4}$; $u_4=7$; $u_5=\dfrac{47}{6}$.
			\item Ta có: $u_n=n+2+\dfrac{5}{n+1}$, do đó $u_n$ nguyên khi và chỉ khi $ \dfrac{5}{n+1}$ nguyên hay $ n+1 $ là ước của 5. Điều đó xảy ra khi $ n+1=5\Leftrightarrow n=4 $. Vậy dãy số có duy nhất một số hạng nguyên là $u_4=7 $.
		\end{enumerate}
		
	}
\end{vd}\dongcham{5}

\begin{vd}%[Đỗ Vũ Minh Thắng]%[1D3B2]
	Cho dãy số $(u_n)$ có số hạng tổng quát $u_n=\dfrac{2n+1}{n+2}$. 
	\begin{tasks}(2)
		\task Viết năm số hạng đầu của dãy số. \dapso{$ 1, \dfrac{5}{4}, \dfrac{7}{5}, \dfrac{3}{2}, \dfrac{11}{7}$}
		\task Tìm số hạng thứ $ 100 $ và $ 200 $. \dapso{$\dfrac{67}{34}$; $ \dfrac{401}{202}$.}
		\task Số $ \dfrac{167}{84}$ là số hạng thứ mấy? \dapso{$ 250 $.}
		\task Dãy số có bao nhiêu số hạng là số nguyên? \dapso{$ 1 $}
	\end{tasks}
	
	\loigiai{
		\begin{enumerate}
			\item Năm số hạng đầu của dãy là:$ u_1=1, u_2=\dfrac{5}{4}, u_3=\dfrac{7}{5}, u_4=\dfrac{3}{2}, u_5=\dfrac{11}{7}$. 
			\item Số hạng thứ 100: $ u_{100}=\dfrac{2.100+1}{100+2}=\dfrac{67}{34}$.\\
			Số hạng thứ 200: $ u_{200}=\dfrac{2.200+1}{200+2}=\dfrac{401}{202}$.
			\item Giả sử $ u_n=\dfrac{167}{84}\Rightarrow \dfrac{2n+1}{n+2}=\dfrac{167}{84}\Leftrightarrow 84(2n+1)=167(n+2)$
			$ \Leftrightarrow n=250 $. \\
			Vậy $ \dfrac{167}{84}$ là số hạng thứ $ 250 $ của dãy số $(u_n)$. 
			\item Ta có: $ u_n=\dfrac{2(n+2)-3}{n+2}=2-\dfrac{3}{n+2}$
			$ \Rightarrow u_n\in \mathbb{Z}\Leftrightarrow \dfrac{3}{n+2}\in \mathbb{Z}\Leftrightarrow 3\;\vdots\; (n+2)\Leftrightarrow n=1 $.\\
			Vậy dãy số có duy nhất một số hạng là số nguyên.
			
		\end{enumerate}
		
	}
\end{vd}\dongcham{10}

\begin{dang}{Tìm các số hạng của dãy số cho bởi công thức truy hồi}
\end{dang}

\begin{vd}
	Cho dãy số $\left(u_n\right)$, biết $\heva{
		& u_1=-1 \\
		& u_{n+1}=u_n+3 \\
	}$ với $n\ge 1$. Tìm ba số hạng đầu tiên của dãy số. \dapso{$-1;2;5.$}
	\loigiai{Ta có $u_1=-1;u_2=u_1+3=2;u_3=u_2+3=5.$}
\end{vd}\dongcham{3}

\begin{vd}
	Cho dãy số $(u_n)$ được xác định bởi$\left\{\begin{aligned}&u_1=1 \\&u_n=2u_{n-1}+1, (n\geq 2).\end{aligned}\right.$ Tính $u_8$.
	\loigiai
	{
		Ta có $u_1=1$, $u_2=3$, $u_3=7$, $u_4=15$, $u_5=31$, $u_6=63$, $u_7=127$, $u_8=255$.
	}
\end{vd}\dongcham{4}

\begin{vd}
	Cho dãy số $(u_n)$ được xác định bởi
	$\left\{\begin{aligned}&u_1=1, u_2=1 \\&u_n=u_{n-1}+u_{n-2}, (n\geq 3)\end{aligned}\right. \text{ (dãy số Phi-bô-na-xi)}.$ Tính $u_7$.
	\loigiai
	{
	Ta có $u_1=1$, $u_2=1$, $u_3=2$, $u_4=3$, $u_5=5$, $u_6=8$, $u_7=13$.
	}
\end{vd}\dongcham{4}


\begin{dang}{Xét sự tăng giảm của dãy số}
	
	\begin{enumerate}[\iconMT]
		\item \indam{Phương pháp 1:} Xét dấu của hiệu số $u_{n+1}-u_n$.
		\begin{itemize}
			\item Nếu $u_{n+1}-u_n>0, \forall n \in \mathbb{N}^*$ thì $(u_n)$ là dãy số tăng.
			\item Nếu $u_{n+1}-u_n<0, \forall n \in \mathbb{N}^*$ thì $(u_n)$ là dãy số giảm.
		\end{itemize}
		\item \indam{Phương pháp 2:} Nếu $u_n>0, \forall n\in \mathbb{N}^*$ thì ta có thể so sánh thương $\dfrac{u_{n+1}}{u_n}$ với $1$.
		\begin{itemize}
			\item Nếu $\dfrac{u_{n+1}}{u_n}>1$ thì $(u_n)$ là dãy số tăng.
			\item Nếu $\dfrac{u_{n+1}}{u_n}<1$ thì $(u_n)$ là dãy số giảm.
		\end{itemize}
		Nếu $u_n<0, \forall n\in \mathbb{N}^*$ thì ta có thể so sánh thương $\dfrac{u_{n+1}}{u_n}$ với $1$.
		\begin{itemize}
			\item Nếu $\dfrac{u_{n+1}}{u_n}<1$ thì $(u_n)$ là dãy số tăng.
			\item Nếu $\dfrac{u_{n+1}}{u_n}>1$ thì $(u_n)$ là dãy số giảm.
		\end{itemize}
	\end{enumerate}
\end{dang}

\begin{vd}
	Xét tính tăng giảm của dãy số sau $(u_n)$ với
	\begin{tasks}(3)
		\task $u_n=n^2$.	\dapso{Tăng.}
		\task $u_n=\dfrac{2n+1}{n+1}$.	\dapso{Tăng.}
		\task $u_n=\dfrac{4^n-1}{4^n+5}$.	\dapso{Tăng.}
		\task $u_n=\dfrac{n}{3^n}$.	\dapso{Giảm.}
		\task $u_n=(-1)^n$.	\dapso{Không tăng, không giảm.}
		\task $u_n=\sqrt{n}-\sqrt{n+3}$.	\dapso{Tăng.}
	\end{tasks}
	\loigiai{
		\begin{enumerate}[a)]
			\item 
			Xét $u_{n+1}-u_n=(n+1)^2-n^2=2n+1 >0, \forall n \in \mathbb{N}^{*}$ nên $(u_n)$ là dãy tăng.
			\item
			Ta có: $u_n=\dfrac{2n+1}{n+1}=2-\dfrac{1}{n+1}$.\\
			$u_{n+1}-u_n=\left(2-\dfrac{1}{n+1+1}\right)-\left(2-\dfrac{1}{n+1}\right)=\dfrac{1}{n+1}-\dfrac{1}{n+2}>0, \forall n \in \mathbb{N}^*$.\\
			Vậy dãy số $(u_n)$ là dãy số tăng.
			\item 
			Ta có $u_n=\dfrac{4^n-1}{4^n+5}=1-\dfrac{6}{4^n+5}$.\\
			$u_{n+1}-u_n=\left(1-\dfrac{6}{4^{n+1}+5}\right)-\left(1-\dfrac{6}{4^n+5}\right)=\dfrac{6}{4^n+5}-\dfrac{6}{4^{n+1}+5} >0, \forall n \in \mathbb{N}^*$.\\
			Vậy dãy số $(u_n)$ là dãy số tăng.
			\item 
			Ta có $u_n=\dfrac{n}{3^n}>0, \forall n \in \mathbb{N}^*$.\\
			Xét thương $\dfrac{u_{n+1}}{u_n}=\dfrac{n+1}{3^{n+1}}:\dfrac{n}{3^n}=\dfrac{n+1}{3.n}<1, \forall  n \in \mathbb{N}^*$.\\
			Vậy $(u_n)$ là dãy số giảm.
			\item
			Ta có:\\ $u_1=(-1)^1=-1.\\
			u_2=(-1)^2=1.\\
			u_3=(-1)^3=-1.$\\
			Vậy $(u_n)$ là dãy không tăng không giảm.
			\item
			Ta có $u_n=\sqrt{n}-\sqrt{n+3}=\dfrac{-3}{\sqrt{n}+\sqrt{n+3}}$.\\
			Xét hiệu\\ 
			$\begin{aligned}
				u_{n+1}-u_n&=\dfrac{-3}{\sqrt{n+1}+\sqrt{n+4}}-\dfrac{-3}{\sqrt{n}+\sqrt{n+3}}\\
				&=\dfrac{3}{\sqrt{n}+\sqrt{n+3}}-\dfrac{3}{\sqrt{n+1}+\sqrt{n+4}}>0, \forall n\in \mathbb{N}^*.
			\end{aligned}$\\
			Vậy $(u_n)$ là dãy số tăng.
		\end{enumerate}
	}
\end{vd}\dongcham{18}


\begin{vd}
	Anh Thanh vừa được tuyển dụng vào một công ty công nghệ, được cam kết lương năm đầu sẽ là $ 200 $ triệu đồng và lương mỗi năm tiếp theo sẽ được tăng thêm $ 25 $ triệu đồng. Gọi $s_n$ (triệu đồng) là lương vào năm thứ $n$ mà anh Thanh làm việc cho công ty đó. Khi đó ta có:
	$$s_1=200, s_n=s_{n-1}+25\, \text{với}\, n \geq 2.$$
	\begin{enumEX}{1}
		\item Tính lương của anh Thanh vào năm thứ $ 5 $ làm việc cho công ty.
		\item Chứng minh $\left(s_n\right)$ là dãy số tăng. Giải thích ý nghĩa thực tế của kết quả này.
	\end{enumEX}
	\loigiai{
		\begin{enumEX}{1}
			\item Ta có: $ s_1=200; s_2=225; s_3=250;s_4=275;s_5=300 $.\\
			Vậy lương của anh Thanh vào năm thứ $ 5 $ là $ 300 $ triệu đồng.
			\item Ta có $ s_n-s_{n-1}=25>0\,\forall n\ge 2 $ do đó $\left(s_n\right)$ là dãy số tăng.\\
			Điều này cho thấy ý nghĩa thực tế là một người có nhiều năm làm việc thì nhận được lương nhiều hơn người mới vào làm việc.
		\end{enumEX}	
	}
\end{vd}\dongcham{6}

\begin{dang}{Xét tính bị chặn của dãy số}
	\begin{itemize}
		\item Để chứng minh dãy số $(u_n)$ bị chặn trên bởi $M$, ta chứng minh $u_n\leq M, \forall n\in \mathbb{N}^*.$
		\item Để chứng minh dãy số $(u_n)$ bị chặn dưới bởi $m$, ta chứng minh $u_n\geq m, \forall n\in \mathbb{N}^*.$
		\item Để chứng minh dãy số bị chặn ta chứng minh nó bị chặn trên và bị chặn dưới.
		\item Nếu dãy số $(u_n)$ tăng thì bị chặn dưới bởi $u_1$; dãy số $(u_n)$ giảm thì bị chặn trên bởi $u_1$.
	\end{itemize}
\end{dang}

\begin{vd}%[1D3G2-4]
	Xét tính bị chặn của các dãy số $(u_n)$ sau, với
	\begin{enumEX}{3}
		\item $u_n=\dfrac{1}{2^n}$.
		\item $u_n=2\sin{n^2}$
		\item $u_n=\dfrac{3n-1}{3n+1}$.
		\item $u_n=\dfrac{n^2+1}{2n^2-3}$.
		\item $u_{n}=\dfrac{1}{n(n+1)}$.
		\item $u_n=\dfrac{2n-1}{\sqrt{n^2+2}}$.
	\end{enumEX}
	\loigiai
	{
		\begin{enumEX}{1}
			\item $0<u_n=\dfrac{1}{2^n}\leq\dfrac{1}{2}$ với mọi $n\in\mathbb{N}^*$ nên dãy $(u_n)$ bị chặn.
			\item Ta có: $-2 \leq u_n=2 \sin n^{2} \leq 2, \forall n \in \mathbb{N}^*$.\\
			Suy ra $(u_n)$ với $u_n=2 \sin n^2$ bị chặn dưới bởi $-2$ và bị chặn trên bời $2$.
			\item Ta có $u_n=\dfrac{3n-1}{3n+1}=1-\dfrac{4}{3n+3}<1, \forall n \geq 1$. Suy ra dãy số $(u_n)$ bị chặn trên.\\
			Mặt khác, $\dfrac{3n-1}{3n+1}>0$, $\forall n\geq 1$ nên dãy số $(u_n)$ bị chặn dưới.\\
			Vậy dãy số $(u_n)$ bị chặn.
			\item Ta có 
			\[u_n=\dfrac{n^2+1}{2n^2-3}=\dfrac{1}{2}+\dfrac{5}{2(2n^2-3)}\tag{1}.\]
			Dễ thấy, với mọi $n \geq 1$ ta có, $-1 \le \dfrac{1}{2n^2-3} \le \dfrac{1}{5}$.\\ Do đó từ $(1)$ suy ra $-2 \le u_n \le 1, \forall n \geq 1$.\\
			Từ đó suy ra $(u_n)$ là một dãy số bị chặn.
			\item Ta có $u_n=\dfrac{1}{n(n+1)}>0, \forall n \geq 1$ nên dãy số bị chặn dưới.\\
			Lại có
			\[u_{n+1}-u_n=\dfrac{1}{(n+1)(n+2)}-\dfrac{1}{n(n+1)}=\dfrac{-2}{n(n+1)(n+2)}<0, \forall n \geq 1.\]
			Vậy dãy số $(u_n)$ là dãy số giảm nên $u_{n+1}\le u_1=\dfrac{1}{2}$. Suy ra dãy số $(u_n)$ bị chặn trên.\\
			Vậy dãy số $(u_n)$ bị chặn.
			\item Ta có $u_n=\dfrac{2n-1}{\sqrt{n^2+2}}>0, \forall n \geq 1$ nên dãy số bị chặn dưới.\\
			Lại có $\dfrac{2n-1}{\sqrt{n^2+2}} \le \dfrac{2n-1}{n} \le \dfrac{2n}{n}=2$ nên dãy số bị chặn trên.\\
			Vậy dãy số $(u_n)$ bị chặn.
		\end{enumEX}	
	}
\end{vd}\dongcham{18}

\begin{vd}%[Hữu Bình]%[1D3B2]
	Chứng minh rằng dãy số $(u_n)$ xác định bởi $u_n=\dfrac{8n+3}{3n+5}$ là một dãy số bị chặn.	\dapso{$ 0<u_n< \dfrac{11}{3}$.}
	\loigiai{
		Ta có $u_n>0, \forall n\geq 1$. Suy ra dãy số bị chặn dưới.\\
		Mặt khác $u_n=\dfrac{8n+3}{3n+5}<\dfrac{8n+3}{3n}=\dfrac{8}{3}+\dfrac{1}{n}<\dfrac{8}{3}+1=\dfrac{11}{3}$.	Do đó dãy số bị chặn trên bởi $\dfrac{11}{3}$.\\
		Vậy dãy số đã cho bị chặn.
	}
\end{vd}\dongcham{4}

\begin{vd}
	Chứng minh rằng dãy số $(u_n)$ với $u_n=\dfrac{3n}{n^2+9}$ bị chặn trên bởi $\dfrac{1}{2}$.
	\loigiai{
		Với mọi $n\geq 1$, ta có $\dfrac{3n}{n^2+9}\leq \dfrac{1}{2} \Leftrightarrow n^2+9\leq 6n \Leftrightarrow (n-3)^2\leq 0$ (đúng).\\
		Vậy dãy số đã cho bị chặn trên bởi $\dfrac{1}{2}$.
	}
\end{vd}\dongcham{4}


\begin{dang}{Vận dụng thực tiễn}
	
\end{dang}

\begin{vd}%[2D2B4-5]%
	Biết rằng năm $ 2016 $, dân số của Việt Nam là
	$ 93{,}422$ triệu người. Tỷ lệ tăng dân số hàng năm của Việt Nam là $ 1{,}07\% $ thì số dân $P_n$ (triệu người) của Việt Nam sau $n$ năm, kể từ năm $2016$, được tính bằng công thức $P_n = 93,422(1 +0,0107 )^n$. Hỏi nếu tăng trưởng theo quy luật như vậy thì vào năm $2026$, dân số Việt Nam khoảng bao nhiêu triệu người?
	\loigiai
	{
		Từ năm 2016 đến năm 2026 là 10 năm, nên $n=10$. Thay vào công thức đã cho ta được
		
		$ P_{10}\approx 1{.}03{,} 9 $ triệu người.
	}
\end{vd}\dongcham{2}

\begin{vd}
	Chị Hương vay trả góp một khoản tiền $ 100 $ triệu đồng và đồng ý trả dần $ 2 $ triệu đồng mỗi tháng với lãi suất $0,8 \%$ số tiền còn lại của mỗi tháng.\\
	Gọi $A_n(n \in \mathbb{N}$) là số tiền còn nợ (triệu đồng) của chị Hương sau $n$ tháng.
	\begin{enumEX}{1}
		\item Tìm lần lượt $A_0, A_1, A_2, A_3, A_4, A_5, A_6$ để tính số tiền còn nợ của chị Hương sau $ 6 $ tháng.
		\item Dự đoán hệ thức truy hồi đối với dãy số $\left(A_n\right)$.
	\end{enumEX}
	\loigiai{
		\begin{enumerate}[a)]
			\item Đặt $A_0=100$ (triệu đồng). Khi đó
			\begin{itemize}
				\item [$\bullet$] $A_1=A_0+\dfrac{0,8}{100}A_0-2=1,008A_0-2=98{,}8$ (triệu đồng).
				\item [$\bullet$] $A_2=A_1+\dfrac{0,8}{100}A_1-2=1,008A_1-2=97,5904$ (triệu đồng).
				\item [$\bullet$] $A_3=A_2+\dfrac{0,8}{100}A_2-2=1,008A_2-2=96,3711232$ (triệu đồng).
				\item [$\bullet$] $A_4=A_3+\dfrac{0,8}{100}A_3-2=1,008A_3-2=95,1420932$ (triệu đồng).
				\item [$\bullet$] $A_5=A_4+\dfrac{0,8}{100}A_4-2=1,008A_4-2=93,9032332$ (triệu đồng).
				\item [$\bullet$] $A_6=A_5+\dfrac{0,8}{100}A_5-2=1,008A_5-2=92,6544632$ (triệu đồng).
			\end{itemize}
			\item Dự đoán hệ thức truy hồi đối với dãy số $\left(A_n\right)$ là $A_n=1,008A_{n-1}-2$, với $n \ge 1$.
		\end{enumerate}
		
	}
\end{vd}\dongcham{8}
\subsection{BÀI TẬP TỰ LUYỆN}
\begin{bt}%[1D3B2-2]%
	Tìm 4 số hạng đầu tiên của các dãy số $(u_n)$ biết số hạng tổng quát:
	\begin{enumEX}{1}
		\item $u_n=\dfrac{1+n}{\sqrt{n^2+1}}$.
		\item $u_n=\dfrac{2^n-1}{2^n+1}$.
		\item $u_n=(-1)^n \sqrt{4^n}$.
		\item $u_n=\dfrac{1+(-1)^n}{n}$.
		\item $u_n=\sin ^2 \dfrac{n \pi}{4}+\cos \dfrac{2n\pi}{3}$.
	\end{enumEX}
	\loigiai{
		\begin{enumEX}{1}
			\item $u_n=\dfrac{1+n}{\sqrt{n^2+1}}$.\\
			+ Với $n=1$ suy ra $u_1=\sqrt{2}$.\\
			+ Với $n=2$ suy ra $u_2=\dfrac{3\sqrt{5}}{5}$.\\
			+ Với $n=3$ suy ra $u_3=\dfrac{2\sqrt{10}}{5}$.\\
			+ Với $n=4$ suy ra $u_4=\dfrac{5\sqrt{17}}{17}$.
			\item $u_n=\dfrac{2^n-1}{2^n+1}$.\\
			+ Với $n=1$ suy ra $u_1=\dfrac{1}{3}$.\\
			+ Với $n=2$ suy ra $u_2=\dfrac{3}{5}$.\\
			+ Với $n=3$ suy ra $u_3=\dfrac{7}{9}$.\\
			+ Với $n=4$ suy ra $u_4=\dfrac{15}{17}$.
			\item $u_n=(-1)^n \sqrt{4^n}$.\\
			+ Với $n=1$ suy ra $u_1=-2$.\\
			+ Với $n=2$ suy ra $u_2=4$.\\
			+ Với $n=3$ suy ra $u_3=-8$.\\
			+ Với $n=4$ suy ra $u_4=16$.
			\item $u_n=\dfrac{1+(-1)^n}{n}$.\\
			+ Với $n=1$ suy ra $u_1=0$.\\
			+ Với $n=2$ suy ra $u_2=1$.\\
			+ Với $n=3$ suy ra $u_3=0$.\\
			+ Với $n=4$ suy ra $u_4=\dfrac{1}{2}$.
			\item $u_n=\sin ^2 \dfrac{n \pi}{4}+\cos \dfrac{2n\pi}{3}$.\\
			+ Với $n=1$ suy ra $u_1=0$.\\
			+ Với $n=2$ suy ra $u_2=\dfrac{1}{2}$.\\
			+ Với $n=3$ suy ra $u_3=\dfrac{3}{2}$.\\
			+ Với $n=4$ suy ra $u_4=-\dfrac{1}{2}$.
		\end{enumEX}
	}
\end{bt}
\cham{10}

\begin{bt}%[1D3B2-2]%
	Cho dãy số $(u_n)$ với số hạng tổng quát $u_n=\dfrac{2n+1}{n+2}$.
	\begin{enumEX}{1}
		\item Tìm $4$ số hạng đầu tiên của dãy số $(u_n)$.
		\item Số $\dfrac{105}{54}$ là số hạng thứ mấy của dãy số $(u_n)$?
	\end{enumEX}
	\loigiai{
		\begin{enumEX}{1}
			\item Tìm $4$ số hạng đầu tiên của dãy số $(u_n)$.\\
			+ Với $n=1$ suy ra $u_1=\dfrac{2\cdot 1+1}{1+2}=1$.\\
			+ Với $n=2$ suy ra $u_1=\dfrac{2\cdot 2+1}{2+2}=\dfrac{5}{4}$.\\
			+ Với $n=3$ suy ra $u_1=\dfrac{2\cdot 3+1}{3+2}=\dfrac{7}{5}$.\\
			+ Với $n=4$ suy ra $u_1=\dfrac{2\cdot 4+1}{4+2}=\dfrac{3}{2}$.
			\item Số $\dfrac{105}{54}$ là số hạng thứ mấy của dãy số $(u_n)$?\\
			Ta có $u_n=\dfrac{105}{54} \Rightarrow \dfrac{2n+1}{n+2}=\dfrac{105}{54}\Rightarrow n=52$.\\
			Vậy số $\dfrac{105}{54}$ là số hạng thứ $52$ của dãy số $(u_n)$.
		\end{enumEX}
	}
\end{bt}
\cham{10}


\begin{bt}
	Viết năm số hạng đầu tiên của dãy số $(u_n)$ và dự đoán công thức số hạng tổng quát $u_n$ theo $n$ của các dãy số $(u_n)$ sau
	\begin{enumEX}{2}
		\item $\left\{\begin{aligned}&u_1=3 \\&u_{n+1}=\sqrt{1+u_n^2}, \forall n\geq 1.\end{aligned}\right.$	\dapso{$ u_n=\sqrt{n+8} $.}
		\item $\left\{\begin{aligned}&u_1=1 \\&u_{n+1}=2u_n+3, \forall n\geq 1.\end{aligned}\right.$	\dapso{$ u_n=2^{n+1}-3 $.}
	\end{enumEX}
	\loigiai
	{
		\begin{enumEX}{1}
			\item Năm số hạng đầu tiên của dãy số đã cho là $u_1=3$, $u_2=\sqrt{10}$, $u_3=\sqrt{11}$, $u_4=\sqrt{12}$, $u_5=\sqrt{13}$.\\
			Ta thấy $u_n>0$ với $n\in\mathbb{N}^*$. Từ công thức $u_{n+1}=\sqrt{1+u_n^2}$ suy ra
			\[u_{n+1}^2=1+u_n^2, \forall n\in \mathbb{N}^* \Leftrightarrow u_{n+1}^2-u_n^2=1, \forall n\in \mathbb{N}^*.\]
			Từ đó ta có
			\allowdisplaybreaks
			\begin{align*}
				u_2^2-u_1^2 &=1\\
				u_3^2-u_2^2 &=1\\
				u_4^2-u_3^2 &=1\\
				\cdots\cdots\cdots\\
				u_n^2-u_{n-1}^2 &=1
			\end{align*}
			Cộng theo vế tất cả các đẳng thức trên ta được
			\[u_n^2-u_1^2= n-1 \Leftrightarrow u_n^2=u_1^2+n-1 \Leftrightarrow u_n^2=n+8 \Leftrightarrow u_n=\sqrt{n+8}.\]
			\item Năm số hạng đầu tiên của dãy số đã cho là $u_1=1$, $u_2=5$, $u_3=13$, $u_4=29$, $u_5=61$.\\
			Từ công thức $u_{n+1}=2u_n+3$ với mọi $n\geq 1$ suy ra
			\[u_{n+1}+3=2u_n+6, \forall n\geq 1 \Leftrightarrow u_{n+1}+3=2(u_n+3), \forall n\geq 1 \Rightarrow \dfrac{u_{n+1}+3}{u_n+3}=2, \forall n\geq 1.\]
			Từ đó ta có
			\allowdisplaybreaks
			\begin{align*}
				\dfrac{u_2+3}{u_1+3} &=2\\
				\dfrac{u_3+3}{u_2+3} &=2\\
				\dfrac{u_4+3}{u_3+3} &=2\\
				\cdots\cdots\\
				\dfrac{u_n+3}{u_{n-1}+3} &=2.\\
			\end{align*}
			Nhân theo vế tất cả các đẳng thức trên ta được
			\[\dfrac{u_n+3}{u_1+3} = 2^{n-1} \Rightarrow u_n+3=2^{n-1}(u_1+3) \Leftrightarrow u_n+3=4\cdot 2^{n-1} \Leftrightarrow u_n=2^{n+1}-3.\]
		\end{enumEX}
	}
\end{bt}\cham{19}


\begin{bt}
	Xét tính tăng giảm của dãy số $(u_n)$, biết
	\begin{tasks}(2)
		\task $u_n=n^3-2n+1$.	\dapso{Tăng.}
		\task $u_n=\dfrac{1}{n}-2$.	\dapso{Giảm.}
		\task $u_n=\dfrac{\sqrt{2}}{3^n}$.	\dapso{Giảm.}
		\task $u_n=\dfrac{1}{n+1}+\dfrac{1}{n+2}+\ldots+\dfrac{1}{2n}$
	\end{tasks}
	\loigiai{
		\begin{enumerate}
			\item Ta có:\\ 
			$\begin{aligned}
				u_{n+1}-u_n&=(n+1)^3-2(n+1)+1-(n^3-2n+1)\\&=n^3+3n^2+3n+1-2n-2-n^3+2n\\
				&=3n-1>0, \forall n \in \mathbb{N}^*.
			\end{aligned}$\\
			Vậy $(u_n)$ là dãy số tăng.
			\item Ta có\\ $u_{n+1}-u_n=\left(\dfrac{1}{n+1}-2\right)-\left(\dfrac{1}{n}-2\right)=\dfrac{1}{n+1}-\dfrac{1}{n}<0, \forall n \in \mathbb{N}^*$.\\
			Vậy $(u_n)$ là dãy số giảm.
			\item Ta có $u_n>0, \forall n \in \mathbb{N}^*$.\\
			Xét thương $\dfrac{u_{n+1}}{u_n}=\dfrac{\sqrt{2}}{3^{n+1}}:\dfrac{\sqrt{2}}{\sqrt{3^2}}=\dfrac{3^{n}}{3^{n+1}}=\dfrac{1}{3}<1$.\\
			Vậy $(u_n)$ là dãy số giảm.
			\item Xét hiệu\\
			$\begin{aligned}
				u_{n+1}-u_n&=\left(\dfrac{1}{n+2}+\dfrac{1}{n+3}+\ldots+\dfrac{1}{2(n+1)}\right)-\left(\dfrac{1}{n+1}+\dfrac{1}{n+2}+\ldots+\dfrac{1}{2n}\right)\\
				&=\dfrac{1}{n+2}-\dfrac{1}{2n+1}-\dfrac{1}{2n+2}\\
				&=\dfrac{1}{2n+2}-\dfrac{1}{2n+1}<0, \forall n\in \mathbb{N}^*.
			\end{aligned}$\\
			Vậy $(u_n)$ là dãy số giảm.
		\end{enumerate}
		
	}
\end{bt}\cham{32}


\begin{bt}%[2D2B4-5]%
	Việt Nam là quốc gia nằm ở phía Đông bán đảo Đông Dương thuộc khu vực Đông Nam Á. Với dân số ước tính $ 93{,}7 $ triệu dân vào đầu năm $ 2018 $, tỉ lệ tăng dân số hàng năm là $ 1{,}2 $\%. Giả sử tỉ lệ tăng dân số từ năm $ 2018 $ đến năm $ 2030 $ không thay đổi thì dân số nước ta đầu năm $ 2030 $ khoảng bao nhiêu?
	\loigiai{
		Từ năm $ 2018 $ đến năm $ 2030 $ là $ 12 $ năm.\\
		Dân số nước ta tính đến năm $ 2030 $ với tỉ lệ tăng dân số không đổi $ 1{,}2 $\% là
		\[
		S=93{,7}\cdot (1+1{,}2\%)^{12}\approx 108{,}12 \text{ triệu dân.}
		\]
	}
\end{bt}\cham{5}

\begin{bt}
	Ông An gửi tiết kiệm $ 100 $ triệu đồng kì hạn $ 1 $ tháng với lãi suất $6 \%$ một năm theo hình thức tính lãi kép. Số tiền (triệu đồng) của ông An thu được sau $n$ tháng được cho bởi công thức
	$$A_n=100\left(1+\dfrac{0,06}{12}\right)^n.$$
	\begin{enumEX}{1}
		\item Tìm số tiền ông An nhận được sau tháng thứ nhất, sau tháng thứ hai.
		\item Tìm số tiền ông An nhận được sau $ 1 $ năm.
	\end{enumEX}
\loigiai{
	\begin{enumerate}[a)]
		\item Số tiền ông An nhận được sau tháng thứ nhất (ứng với $n = 1$) là:
		$$ A_1 = 100\left(1 + \dfrac{0{,}06}{12}\right)^1 = 100(1 + 0{,}005) = 100{,}5 \text{ (triệu đồng)}. $$
		Số tiền ông An nhận được sau tháng thứ hai (ứng với $n = 2$) là:
		$$ A_2 = 100\left(1 + \dfrac{0{,}06}{12}\right)^2 = 100(1{,}005)^2 = 101{,}0025 \text{ (triệu đồng)}. $$
		
		\item Một năm có $12$ tháng. Số tiền ông An nhận được sau $1$ năm (ứng với $n = 12$) là:
		$$ A_{12} = 100\left(1 + \dfrac{0{,}06}{12}\right)^{12} = 100(1{,}005)^{12} \approx 106{,}168 \text{ (triệu đồng)}. $$
	\end{enumerate}
}
\end{bt}\cham{8}


\begin{bt}
	Chị Mai gửi tiền tiết kiệm vào ngân hàng theo thể thức lãi kép như sau: Lần đầu chị gửi $ 100 $ triệu đồng. Sau đó, cứ hết $ 1 $ tháng chị lại gửi thêm vào ngân hàng $ 6 $ triệu đồng. Biết lãi suất của ngân hàng là $0,5 \%$ một tháng. Gọi $P_n$ (triệu đồng) là số tiền chị có trong ngân hàng sau $n$ tháng.
	\begin{enumEX}{1}
		\item Tính số tiền chị có trong ngân hàng sau $ 1 $ tháng.
		\item Tính số tiền chị có trong ngân hàng sau $ 3 $ tháng.
		\item Dự đoán hệ thức truy hồi đối với dãy $P_n$.
	\end{enumEX}
	\loigiai{
		Đặt $P_0=100$ (triệu đồng). Ta có
		\begin{enumerate}[a)]
			\item Số tiền của chị Mai sau 1 tháng là $P_1=P_0 + P_0 \cdot \dfrac{0,5}{100}+6=106,5$ (triệu đồng)
			\item Số tiền của chị Mai sau 2 tháng là $P_2=P_1 + P_1 \cdot \dfrac{0,5}{100}+6=113,0325$ (triệu đồng).\\
			Số tiền của chị Mai sau 3 tháng là $P_3=P_2 + P_2 \cdot \dfrac{0,5}{100}+6=119,598$ (triệu đồng).
			\item Dự đoán hệ thức truy hồi đối với dãy $P_n$ là $P_n=1{,}005 \cdot P_{n-1}+6.$
	\end{enumerate}}
\end{bt}\cham{10}



